\section{Computational complexity}
\label{sec:complexity}

Asymptotic cost decides which methods are \emph{admissible} at a given sample size, dimension and simulation cost,
before accuracy is considered. Table~\ref{tab:complexity} collects fit, storage and prediction costs for the families
reviewed above, in the usual algorithm for each; Table~\ref{tab:complexity_upscale} does the same for the
representations used to make a learner optimizer-ready. Constants matter at engineering sample sizes, so the study
also reports measured times.

Three walls recur. Exact GPs and dense RBF systems need $O(n^3)$ time and $O(n^2)$ memory
\cite{rasmussen2006gaussian,wendland2004scattered}---about 0.8\,GB of kernel matrix at $n=10^4$---which inducing-point,
Nystr\"om and random-feature approximations reduce to $O(nq^2)$ \cite{titsias2009variational,williams2001using,rahimi2007random},
and partitioned or local GPs reduce by working on subsets \cite{gramacy2008bayesian,gramacy2015local}. Full tensor
grids, and hence tensor-product spline upscaling, cost $O(m^d)$: 24 points per axis is $1.4\times10^4$ nodes in three
dimensions, $1.9\times10^8$ in six and $6\times10^{13}$ in ten; sparse grids \cite{bungartz2004sparse} and tensor-train
formats \cite{oseledets2011tensor,oseledets2010tt} restore near-linear scaling in $d$ under mixed smoothness or low rank.
Delaunay-based interpolants (piecewise-linear, natural neighbour) grow as $O(n^{\lceil d/2\rceil})$ in the worst case and
are impractical beyond about six inputs. By contrast, tree ensembles fit in $O(n\log n)$ per tree and predict in time
independent of $n$ \cite{breiman2001random,chen2016xgboost,ke2017lightgbm}, and in-context tabular models replace training
by $O(n^2)$ attention \cite{hollmann2025tabpfn}.

Loops multiply these costs. An active-learning loop that refits an exact GP after each of $N$ acquisitions costs
$O(N^4)$ unless the Cholesky factor is updated, and joint co-kriging over $L$ fidelities costs $O((\sum_\ell n_\ell)^3)$ while
the recursive formulation costs $O(\sum_\ell n_\ell^3)$ \cite{legratiet2014recursive}. Randomised linear algebra
\cite{halko2011finding,erichson2019randomized} reduces the cost of a decomposition to its intrinsic rank, which is how
compressed DMD accelerates its intermediate step. For finite-element simulators costing an hour per run, every surrogate
term is negligible up to $n\sim10^4$; for cheap simulators these terms decide the winner.

\begin{table}[ht]
\centering\small
\caption{Asymptotic costs of surrogate learners. $n$ training points, $d$ inputs, $P$ basis terms, $s$ active terms,
$q$ inducing points or features, $B$ trees, $T$ boosting rounds, $D$ depth, $W$ weights, $E$ epochs, $K$ ensemble
members, $H R$ hyperparameter iterations $\times$ restarts.}
\label{tab:complexity}
\resizebox{\linewidth}{!}{%
\begin{tabular}{@{}llll@{}}
\toprule
Method & Fit time & Memory & Predict (mean / variance) \\
\midrule
Polynomial RSM, $P=\binom{d+p}{p}$ & $O(nP^2+P^3)$ & $O(nP)$ & $O(P)$ / $O(P^2)$ \\
Sparse regression, sparse PCE (OMP) & $O(snP)$ & $O(nP)$ & $O(s)$ \\
Exact GP / kriging & $O(HR\,n^3)$ & $O(n^2)$ & $O(nd)$ / $O(n^2)$ \\
Sparse GP, Nystr\"om, random features & $O(nq^2)$ & $O(nq)$ & $O(q)$ / $O(q^2)$ \\
Treed GP (leaves $n_j$) & $O(S\sum_j n_j^3)$ & $O(\sum_j n_j^2)$ & $O(n_j)$ \\
RBF interpolation (dense) & $O(n^3)$ & $O(n^2)$ & $O(nd)$ \\
IDW / $k$-NN & $O(dn\log n)$ (tree build) & $O(nd)$ & $O(\log n)$--$O(nd)$ \\
Delaunay linear / natural neighbour & $O(n^{\lceil d/2\rceil})$ worst case & same & $O(\log n)$ \\
MARS, $M$ terms & $O(ndM^2)$ & $O(nM)$ & $O(M)$ \\
Random forest & $O(Bd'n\log n)$ & $O(Bn)$ & $O(BD)$ \\
Gradient boosting (histogram) & $O(Tnd)$ & $O(nd+T2^D)$ & $O(TD)$ \\
BART, $m_t$ trees, $S$ sweeps & $O(Sm_tn)$ & $O(m_tn)$ & $O(Sm_tD)$ \\
MLP / deep ensemble & $O(KEnW)$ & $O(KW)$ & $O(KW)$ \\
Deep kernel learning & $O(E(n^3+nW))$ & $O(n^2+W)$ & $O(nd+W)$ / $O(n^2)$ \\
TabPFN (in-context) & --- & $O(n^2)$ & batched $O(nM)$ \\
Tensor train (cross) & $O(dmr^3)$; $O(dmr^2)$ samples & $O(dmr^2)$ & $O(dr^2)$ \\
Compressed sensing (OMP), basis $N$ & $O(snN)$ & $O(nN)$ & $O(s)$ \\
\bottomrule
\end{tabular}}
\end{table}

\begin{table}[ht]
\centering\small
\caption{Costs of optimizer-ready representations. $m$ nodes per axis, spline order $k$ ($k=4$ cubic), Smolyak level
$\lambda$ with $N_\lambda=O(2^\lambda\lambda^{d-1})$ nodes, rank $r$, $K$ local models fitted to $n_{\mathrm{fit}}$
model predictions, $c_M$ the cost of one call to the intermediate model ($O(nd)$ for a GP mean, $O(n^2)$ with its
variance). Gradients of the interpolants are analytic; by finite differences they cost $2d$ evaluations.}
\label{tab:complexity_upscale}
\resizebox{\linewidth}{!}{%
\begin{tabular}{@{}lllll@{}}
\toprule
Representation & Build & Storage & Evaluate & Gradient \\
\midrule
Intermediate model itself (no upscaling) & --- & model's & $c_M$ & $O(c_M)$ analytic; $2d\,c_M$ by differences \\
Tensor B-spline on dense grid & $m^d c_M + O(dm^d)$ & $O(m^d)$ & $O(k^d)$ & $O(dk^d)$ \\
Multilevel B-spline (MBA) & $O(h(n+m_h^d))$ & $O(m_h^d)$ & $O(k^d)$ & $O(dk^d)$ \\
Smolyak sparse grid, global (barycentric) basis & $N_\lambda c_M$ & $O(N_\lambda)$ & $O(N_\lambda)$ & $O(dN_\lambda)$ \\
Additive P-spline / GAM & $O(ndK)$ & $O(dK)$ & $O(dk)$ & $O(dk)$ \\
Tensor-train grid, local / global axis basis & $O(dmr^2)\,c_M + O(dmr^3)$ & $O(dmr^2)$ & $O(dr^2)$ / $O(dmr^2)$
  & $O(d^2r^2)$ / $O(d^2mr^2)$ \\
Partition of unity of $K$ local quadratics & $n_{\mathrm{fit}}c_M + O(Kn_{\mathrm{fit}}d^4)$ & $O(Kd^2)$ & $O(Kd^2)$
  & $O(Kd^2)$ \\
\bottomrule
\end{tabular}}
\end{table}

\paragraph{When does upscaling pay for itself?}
Upscaling never reduces the cost of \emph{training} the intermediate model (an exact GP stays $O(n^3)$); it trades a
one-off build for cheaper evaluations. Let a representation need $M$ calls to the intermediate model
(Table~\ref{tab:complexity_upscale}: $M=m^d$, $N_\lambda$, $O(dmr^2)$ or $n_{\mathrm{fit}}$) plus fitting work
$C_{\mathrm{fit}}$, and let one evaluation cost $c_U$ after upscaling and $c_M$ before. Over $N$ optimizer evaluations
the build is repaid when
\begin{equation}
  N\,(c_M-c_U) \;>\; M\,c_M + C_{\mathrm{fit}}
  \quad\Longleftrightarrow\quad
  N \;>\; \frac{M\,c_M + C_{\mathrm{fit}}}{c_M-c_U} \;\approx\; M
  \qquad (c_U \ll c_M,\ C_{\mathrm{fit}} \ll M c_M).
  \label{eq:upscale_breakeven}
\end{equation}
Three consequences follow. (i) Because the build consists of model calls, the break-even budget is almost independent
of $n$: a larger GP makes the build dearer by the same factor by which it makes each saved evaluation dearer, while the
per-evaluation speed-up $c_M/c_U$ grows as $n$ for the GP mean and as $n^2$ when the variance is needed (expected
improvement and other uncertainty-driven searches). (ii) The dimension decides admissibility through $M$: a tensor grid
with $m=12$ has $M=144$ nodes at $d=2$ but $2.5\times10^5$ at $d=5$, so a typical optimizer budget repays it only in low
dimension; tensor-train and local-quadratic representations keep $M$ polynomial in $d$, the former only when the
function is low-rank. (iii) A representation whose evaluation scales with its own size, such as a Smolyak grid evaluated
through a global basis ($c_U=O(N_\lambda)$), is faster than the intermediate model only when $N_\lambda<n$---never for
a GP trained on fewer points than the grid has nodes.

The E8 timings (GP intermediate, $n\le500$) follow these orders. GP evaluation grew linearly in $n$ (1.8, 2.8, 9.1 and
13.5\,\textmu s per point at $n=40$, 60, 200 and 300 in two dimensions). Tensor-spline evaluation did not depend on $n$
(0.5\,\textmu s) and rose with $d$ (18\,\textmu s at $d=5$: a 36-fold rise against the $4^d$ ratio of 64). The $d=5$,
$m=12$ build decomposed into $Mc_M$ plus an $n$-independent fit (11.8\,s $\approx$ 1.1\,s $+$ 10.7\,s at $n=100$;
19.2\,s $\approx$ 7.5\,s $+$ 11.7\,s at $n=500$). With $N\sim10^4$ optimizer evaluations per task the spline paid at
once in two dimensions ($M=144$--$576$; optimizer time 5.5\,s $\to$ 1.6\,s with no loss of accuracy) but not in five
($M\approx2.5\times10^5$). The Smolyak grid ($N_\lambda\le2\times10^4$) was slower than every GP it upscaled
(180--500\,\textmu s per point). For the small GPs of E8 the savings are seconds. The regime in which GP upscaling
matters---$n$ of $10^4$--$10^5$, where $c_M$ reaches milliseconds and $c_M/c_U$ exceeds $10^3$---lies beyond the study's
design.
